
\documentclass{article}
\usepackage{colortbl}
\usepackage{makecell}
\usepackage{multirow}
\usepackage{supertabular}

\begin{document}

\newcounter{utterance}

\centering \large Interaction Transcript for game `clean\_up', experiment `0\_easy\_7obj\_es', episode 2 with deepseek{-}v4{-}pro.
\vspace{24pt}

{ \footnotesize  \setcounter{utterance}{1}
\setlength{\tabcolsep}{0pt}
\begin{supertabular}{c@{$\;$}|p{.15\linewidth}@{}p{.15\linewidth}p{.15\linewidth}p{.15\linewidth}p{.15\linewidth}p{.15\linewidth}}
    \# & $\;$A & \multicolumn{4}{c}{Game Master} & $\;\:$B\\
    \hline

    \theutterance \stepcounter{utterance}  
    & & \multicolumn{4}{p{0.6\linewidth}}{
        \cellcolor[rgb]{0.9,0.9,0.9}{
            \makecell[{{p{\linewidth}}}]{
                \texttt{\tiny{[P1$\langle$GM]}}
                \texttt{Soy tu maestro de juego, y estás jugando un juego colaborativo con la siguiente cuadrícula como tablero de juego:} \\
\\ 
\texttt{```} \\
\texttt{ 1234567} \\
\texttt{╔═══════╗} \\
\texttt{║◌T◌◌◌◌◌║ 1} \\
\texttt{║A◌◌◌◌◌D║ 2} \\
\texttt{║◌◌◌◌◌┌─╢ 3} \\
\texttt{╟──┐O◌│◌║ 4} \\
\texttt{║◌P├──┼─╢ 5} \\
\texttt{║◌S│M◌│◌║ 6} \\
\texttt{║◌◌│◌◌│◌║ 7} \\
\texttt{╚══╧══╧═╝} \\
\\ 
\texttt{```} \\
\\ 
\texttt{* El borde superior muestra coordenadas x que aumentan hacia la derecha, y el borde derecho coordenadas y que aumentan hacia abajo.} \\
\texttt{* Los siguientes objetos están colocados aleatoriamente en tu cuadrícula: 'P', 'O', 'T', 'S', 'D', 'A', 'M'.} \\
\\ 
\\ 
\texttt{El otro jugador ve una variación del tablero de juego, donde los objetos están colocados en diferentes ubicaciones aleatorias. No puedes ver el tablero del otro jugador, y él no puede ver el tuyo.} \\
\\ 
\\ 
\texttt{**Objetivo:**} \\
\\ 
\texttt{Ambos jugadores necesitan mover los objetos en sus respectivos fondos de manera que los objetos idénticos terminen en las mismas coordenadas. Tienes que comunicarte con el otro jugador para acordar una configuración objetivo común.} \\
\\ 
\\ 
\texttt{**Reglas:**} \\
\\ 
\texttt{* En cada turno, puedes enviar exactamente uno de los dos comandos siguientes:} \\
\texttt{ 1. `DECIR: <MENSAJE>`: para enviar un mensaje (todo hasta el siguiente salto de línea) al otro jugador. Se lo reenviaré a tu compañero.} \\
\texttt{2. `MOVERR: <OBJETO>, (<X>, <Y>)`: para mover un objeto a una nueva posición, donde `<X>` es la columna y `<Y>` es la fila. Te informaré si tu movimiento fue exitoso o no.} \\
\texttt{* Si no sigues el formato, o envías varios comandos a la vez, tendré que penalizarte.} \\
\texttt{* Si ambos jugadores acumulan más de 16 penalizaciones, ambos pierden el juego.} \\
\texttt{* ¡Es vital que te comuniques con el otro jugador respecto a tu estado objetivo! ¡La *única* forma de transmitir tu estrategia al otro jugador es usando el comando `DECIR: <MENSAJE>`!} \\
\\ 
\\ 
\texttt{**Mover objetos**} \\
\\ 
\texttt{* Solo puedes mover objetos a celdas dentro de los límites de la cuadrícula. La celda de destino debe estar vacía, es decir, solo debe contener el símbolo '◌'.} \\
\texttt{* Si intentas mover un objeto a un lugar que no está vacío, o intentas moverlo fuera de la cuadrícula, tendré que penalizarte. Tendrás otro intento.} \\
\texttt{* ¡Antes de hacer un movimiento, verifica dos veces que el lugar de destino esté vacío, y no contenga ninguna letra, marco o línea!} \\
\\ 
\\ 
\texttt{**Fin del juego**} \\
\\ 
\texttt{Si crees que alcanzaste el objetivo de alinear todos los objetos, puedes pedirle al otro jugador que termine el juego enviando `DECIR: ¿terminado?`. Si el otro jugador te pide que termines el juego, y respondes con `DECIR: ¡terminado!`, el juego terminará.} \\
\\ 
\texttt{Ambos jugadores ganan si el juego termina dentro de 28 rondas, donde una ronda se define como dos jugadores enviando cada uno un comando válido.} \\
\\ 
\texttt{**Puntuación:**} \\
\\ 
\texttt{Cuanto más cerca estén los objetos idénticos en ambos tableros de juego, más puntos obtendrás. Las penalizaciones reducen tus puntos. ¿Puedes batir el récord?} \\
\\ 
\\ 
\texttt{¡Por favor, envía un mensaje al otro jugador para iniciar el juego!} \\
            }
        }
    }
    & \\ \\

    \theutterance \stepcounter{utterance}  
    & \multicolumn{4}{p{0.6\linewidth}}{
        \cellcolor[rgb]{0.9,0.9,0.9}{
            \makecell[{{p{\linewidth}}}]{
                \texttt{\tiny{[P1$\rangle$GM]}}
                \texttt{Parece que mi personaje está basado en el tablero que me mostraste. ¡Qué intriga! Veo que la cuadrícula tiene varios objetos dispersos y necesito coordinar sus ubicaciones con mi compañero.} \\
\\ 
\texttt{Para comenzar, haré un reconocimiento inicial. En mi tablero veo las siguientes posiciones:} \\
\\ 
\texttt{{-} **P** en columna 2, fila 5 (dentro de un cuadro en la parte inferior)} \\
\texttt{{-} **O** en columna 3, fila 4 (dentro de un área con líneas)} \\
\texttt{{-} **T** en columna 2, fila 1} \\
\texttt{{-} **S** en columna 2, fila 6} \\
\texttt{{-} **D** en columna 7, fila 2} \\
\texttt{{-} **A** en columna 1, fila 2} \\
\texttt{{-} **M** en columna 3, fila 6} \\
\\ 
\texttt{También veo que hay paredes y estructuras que podrían bloquear el movimiento, pero los objetos están dentro de celdas accesibles.} \\
\\ 
\texttt{Para iniciar la coordinación, enviaré un primer mensaje a mi compañero:} \\
\\ 
\texttt{{-}{-}{-}} \\
\\ 
\texttt{**DECIR: Hola, ¿puedes indicarme las coordenadas de tus objetos P, O, T, S, D, A, M? Por mi parte, estas son mis posiciones actuales: P(2,5), O(3,4), T(2,1), S(2,6), D(7,2), A(1,2), M(3,6). ¿Cuáles tienes tú?**} \\
            }
        }
    }
    & & \\ \\

    \theutterance \stepcounter{utterance}  
    & & & \multicolumn{2}{p{0.3\linewidth}}{
        \cellcolor[rgb]{0.9,0.9,0.9}{
            \makecell[{{p{\linewidth}}}]{
                \texttt{\tiny{[GM$|$GM]}}
                \texttt{Parece que mi personaje está basado en el tablero que me mostraste. ¡Qué intriga! Veo que la cuadrícula tiene varios objetos dispersos y necesito coordinar sus ubicaciones con mi compañero.} \\
\\ 
\texttt{Para comenzar, haré un reconocimiento inicial. En mi tablero veo las siguientes posiciones:} \\
\\ 
\texttt{{-} **P** en columna 2, fila 5 (dentro de un cuadro en la parte inferior)} \\
\texttt{{-} **O** en columna 3, fila 4 (dentro de un área con líneas)} \\
\texttt{{-} **T** en columna 2, fila 1} \\
\texttt{{-} **S** en columna 2, fila 6} \\
\texttt{{-} **D** en columna 7, fila 2} \\
\texttt{{-} **A** en columna 1, fila 2} \\
\texttt{{-} **M** en columna 3, fila 6} \\
\\ 
\texttt{También veo que hay paredes y estructuras que podrían bloquear el movimiento, pero los objetos están dentro de celdas accesibles.} \\
\\ 
\texttt{Para iniciar la coordinación, enviaré un primer mensaje a mi compañero:} \\
\\ 
\texttt{{-}{-}{-}} \\
\\ 
\texttt{**DECIR: Hola, ¿puedes indicarme las coordenadas de tus objetos P, O, T, S, D, A, M? Por mi parte, estas son mis posiciones actuales: P(2,5), O(3,4), T(2,1), S(2,6), D(7,2), A(1,2), M(3,6). ¿Cuáles tienes tú?**} \\
            }
        }
    }
    & & \\ \\

    \theutterance \stepcounter{utterance}  
    & & & \multicolumn{2}{p{0.3\linewidth}}{
        \cellcolor[rgb]{0.9,0.9,0.9}{
            \makecell[{{p{\linewidth}}}]{
                \texttt{\tiny{[GM$|$GM]}}
                \texttt{Invalid format: head is not empty: 'Parece que mi personaje está basado en el tablero que me mostraste. ¡Qué intriga! Veo que la cuadrícula tiene varios objetos dispersos y necesito coordinar sus ubicaciones con mi compañero.} \\
\\ 
\texttt{Para comenzar, haré un reconocimiento inicial. En mi tablero veo las siguientes posiciones:} \\
\\ 
\texttt{{-} **P** en columna 2, fila 5 (dentro de un cuadro en la parte inferior)} \\
\texttt{{-} **O** en columna 3, fila 4 (dentro de un área con líneas)} \\
\texttt{{-} **T** en columna 2, fila 1} \\
\texttt{{-} **S** en columna 2, fila 6} \\
\texttt{{-} **D** en columna 7, fila 2} \\
\texttt{{-} **A** en columna 1, fila 2} \\
\texttt{{-} **M** en columna 3, fila 6} \\
\\ 
\texttt{También veo que hay paredes y estructuras que podrían bloquear el movimiento, pero los objetos están dentro de celdas accesibles.} \\
\\ 
\texttt{Para iniciar la coordinación, enviaré un primer mensaje a mi compañero:} \\
\\ 
\texttt{{-}{-}{-}} \\
\\ 
\texttt{**'} \\
            }
        }
    }
    & & \\ \\

    \theutterance \stepcounter{utterance}  
    & & \multicolumn{4}{p{0.6\linewidth}}{
        \cellcolor[rgb]{0.9,0.9,0.9}{
            \makecell[{{p{\linewidth}}}]{
                \texttt{\tiny{[P1$\langle$GM]}}
                \texttt{Penalización: ¡Tu mensaje no debe contener nada antes del comando!} \\
\\ 
\texttt{Asegúrate de que tu respuesta solo contenga `DECIR: <MENSAJE>` o `MOVERR: <OBJETO>, (<X>, <Y>)`, ¡y nada más!} \\
\\ 
\texttt{Han acumulado colectivamente 1 de 16 penalizaciones.} \\
\texttt{¡Por favor, inténtalo de nuevo!} \\
            }
        }
    }
    & \\ \\

    \theutterance \stepcounter{utterance}  
    & \multicolumn{4}{p{0.6\linewidth}}{
        \cellcolor[rgb]{0.9,0.9,0.9}{
            \makecell[{{p{\linewidth}}}]{
                \texttt{\tiny{[P1$\rangle$GM]}}
                \texttt{DECIR: Hola, mis posiciones actuales son P(2,5), O(3,4), T(2,1), S(2,6), D(7,2), A(1,2), M(3,6). Por favor dime las tuyas para coordinarnos.} \\
            }
        }
    }
    & & \\ \\

    \theutterance \stepcounter{utterance}  
    & & & \multicolumn{2}{p{0.3\linewidth}}{
        \cellcolor[rgb]{0.9,0.9,0.9}{
            \makecell[{{p{\linewidth}}}]{
                \texttt{\tiny{[GM$|$GM]}}
                \texttt{DECIR: Hola, mis posiciones actuales son P(2,5), O(3,4), T(2,1), S(2,6), D(7,2), A(1,2), M(3,6). Por favor dime las tuyas para coordinarnos.} \\
            }
        }
    }
    & & \\ \\

    \theutterance \stepcounter{utterance}  
    & & \multicolumn{4}{p{0.6\linewidth}}{
        \cellcolor[rgb]{0.9,0.9,0.9}{
            \makecell[{{p{\linewidth}}}]{
                \texttt{\tiny{[GM$\rangle$P2]}}
                \texttt{Soy tu maestro de juego, y estás jugando un juego colaborativo con la siguiente cuadrícula como tablero de juego:} \\
\\ 
\texttt{```} \\
\texttt{ 1234567} \\
\texttt{╔═══════╗} \\
\texttt{║◌◌P◌A◌O║ 1} \\
\texttt{║S◌◌◌M◌◌║ 2} \\
\texttt{║◌◌◌◌D┌─╢ 3} \\
\texttt{╟──┐◌◌│◌║ 4} \\
\texttt{║◌◌├──┼─╢ 5} \\
\texttt{║◌◌│◌◌│T║ 6} \\
\texttt{║◌◌│◌◌│◌║ 7} \\
\texttt{╚══╧══╧═╝} \\
\\ 
\texttt{```} \\
\\ 
\texttt{* El borde superior muestra coordenadas x que aumentan hacia la derecha, y el borde derecho coordenadas y que aumentan hacia abajo.} \\
\texttt{* Los siguientes objetos están colocados aleatoriamente en tu cuadrícula: 'P', 'O', 'T', 'S', 'D', 'A', 'M'.} \\
\\ 
\\ 
\texttt{El otro jugador ve una variación del tablero de juego, donde los objetos están colocados en diferentes ubicaciones aleatorias. No puedes ver el tablero del otro jugador, y él no puede ver el tuyo.} \\
\\ 
\\ 
\texttt{**Objetivo:**} \\
\\ 
\texttt{Ambos jugadores necesitan mover los objetos en sus respectivos fondos de manera que los objetos idénticos terminen en las mismas coordenadas. Tienes que comunicarte con el otro jugador para acordar una configuración objetivo común.} \\
\\ 
\\ 
\texttt{**Reglas:**} \\
\\ 
\texttt{* En cada turno, puedes enviar exactamente uno de los dos comandos siguientes:} \\
\texttt{ 1. `DECIR: <MENSAJE>`: para enviar un mensaje (todo hasta el siguiente salto de línea) al otro jugador. Se lo reenviaré a tu compañero.} \\
\texttt{2. `MOVERR: <OBJETO>, (<X>, <Y>)`: para mover un objeto a una nueva posición, donde `<X>` es la columna y `<Y>` es la fila. Te informaré si tu movimiento fue exitoso o no.} \\
\texttt{* Si no sigues el formato, o envías varios comandos a la vez, tendré que penalizarte.} \\
\texttt{* Si ambos jugadores acumulan más de 16 penalizaciones, ambos pierden el juego.} \\
\texttt{* ¡Es vital que te comuniques con el otro jugador respecto a tu estado objetivo! ¡La *única* forma de transmitir tu estrategia al otro jugador es usando el comando `DECIR: <MENSAJE>`!} \\
\\ 
\\ 
\texttt{**Mover objetos**} \\
\\ 
\texttt{* Solo puedes mover objetos a celdas dentro de los límites de la cuadrícula. La celda de destino debe estar vacía, es decir, solo debe contener el símbolo '◌'.} \\
\texttt{* Si intentas mover un objeto a un lugar que no está vacío, o intentas moverlo fuera de la cuadrícula, tendré que penalizarte. Tendrás otro intento.} \\
\texttt{* ¡Antes de hacer un movimiento, verifica dos veces que el lugar de destino esté vacío, y no contenga ninguna letra, marco o línea!} \\
\\ 
\\ 
\texttt{**Fin del juego**} \\
\\ 
\texttt{Si crees que alcanzaste el objetivo de alinear todos los objetos, puedes pedirle al otro jugador que termine el juego enviando `DECIR: ¿terminado?`. Si el otro jugador te pide que termines el juego, y respondes con `DECIR: ¡terminado!`, el juego terminará.} \\
\\ 
\texttt{Ambos jugadores ganan si el juego termina dentro de 28 rondas, donde una ronda se define como dos jugadores enviando cada uno un comando válido.} \\
\\ 
\texttt{**Puntuación:**} \\
\\ 
\texttt{Cuanto más cerca estén los objetos idénticos en ambos tableros de juego, más puntos obtendrás. Las penalizaciones reducen tus puntos. ¿Puedes batir el récord?} \\
\\ 
\texttt{El otro jugador inició el juego enviando este mensaje:} \\
\\ 
\texttt{"Hola, mis posiciones actuales son P(2,5), O(3,4), T(2,1), S(2,6), D(7,2), A(1,2), M(3,6). Por favor dime las tuyas para coordinarnos."} \\
\\ 
\texttt{¿Cuál es tu primer comando?} \\
            }
        }
    }
    & \\ \\

    \theutterance \stepcounter{utterance}  
    & & & \multicolumn{4}{p{0.6\linewidth}}{
        \cellcolor[rgb]{0.9,0.9,0.9}{
            \makecell[{{p{\linewidth}}}]{
                \texttt{\tiny{[GM$\langle$P2]}}
                \texttt{Mi primer comando será para compartir mis posiciones actuales:} \\
\\ 
\texttt{```} \\
\texttt{DECIR: Hola, mis posiciones actuales son P(3,1), O(7,1), T(7,6), S(1,2), D(5,3), A(5,1), M(5,2)} \\
\texttt{```} \\
            }
        }
    }
     \\ \\

    \theutterance \stepcounter{utterance}  
    & & & \multicolumn{2}{p{0.3\linewidth}}{
        \cellcolor[rgb]{0.9,0.9,0.9}{
            \makecell[{{p{\linewidth}}}]{
                \texttt{\tiny{[GM$|$GM]}}
                \texttt{Mi primer comando será para compartir mis posiciones actuales:} \\
\\ 
\texttt{```} \\
\texttt{DECIR: Hola, mis posiciones actuales son P(3,1), O(7,1), T(7,6), S(1,2), D(5,3), A(5,1), M(5,2)} \\
\texttt{```} \\
            }
        }
    }
    & & \\ \\

    \theutterance \stepcounter{utterance}  
    & & & \multicolumn{2}{p{0.3\linewidth}}{
        \cellcolor[rgb]{0.9,0.9,0.9}{
            \makecell[{{p{\linewidth}}}]{
                \texttt{\tiny{[GM$|$GM]}}
                \texttt{Invalid format: head is not empty: 'Mi primer comando será para compartir mis posiciones actuales:} \\
\\ 
\\ 
\texttt{'} \\
            }
        }
    }
    & & \\ \\

    \theutterance \stepcounter{utterance}  
    & & \multicolumn{4}{p{0.6\linewidth}}{
        \cellcolor[rgb]{0.9,0.9,0.9}{
            \makecell[{{p{\linewidth}}}]{
                \texttt{\tiny{[GM$\rangle$P2]}}
                \texttt{Penalización: ¡Tu mensaje no debe contener nada antes del comando!} \\
\\ 
\texttt{Asegúrate de que tu respuesta solo contenga `DECIR: <MENSAJE>` o `MOVERR: <OBJETO>, (<X>, <Y>)`, ¡y nada más!} \\
\\ 
\texttt{Han acumulado colectivamente 2 de 16 penalizaciones.} \\
\texttt{¡Por favor, inténtalo de nuevo!} \\
            }
        }
    }
    & \\ \\

    \theutterance \stepcounter{utterance}  
    & & & \multicolumn{4}{p{0.6\linewidth}}{
        \cellcolor[rgb]{0.9,0.9,0.9}{
            \makecell[{{p{\linewidth}}}]{
                \texttt{\tiny{[GM$\langle$P2]}}
                \texttt{DECIR: Hola, mis posiciones actuales son P(3,1), O(7,1), T(7,6), S(1,2), D(5,3), A(5,1), M(5,2)} \\
            }
        }
    }
     \\ \\

    \theutterance \stepcounter{utterance}  
    & & & \multicolumn{2}{p{0.3\linewidth}}{
        \cellcolor[rgb]{0.9,0.9,0.9}{
            \makecell[{{p{\linewidth}}}]{
                \texttt{\tiny{[GM$|$GM]}}
                \texttt{DECIR: Hola, mis posiciones actuales son P(3,1), O(7,1), T(7,6), S(1,2), D(5,3), A(5,1), M(5,2)} \\
            }
        }
    }
    & & \\ \\

    \theutterance \stepcounter{utterance}  
    & & \multicolumn{4}{p{0.6\linewidth}}{
        \cellcolor[rgb]{0.9,0.9,0.9}{
            \makecell[{{p{\linewidth}}}]{
                \texttt{\tiny{[P1$\langle$GM]}}
                \texttt{Tu mensaje ha sido transmitido al otro jugador.} \\
\\ 
\\ 
\texttt{Actualmente estás jugando la ronda 1 de un máximo de 28.} \\
\texttt{Han acumulado colectivamente 2 de 16 penalizaciones.} \\
\texttt{El otro jugador envió este mensaje:} \\
\\ 
\\ 
\texttt{"Hola, mis posiciones actuales son P(3,1), O(7,1), T(7,6), S(1,2), D(5,3), A(5,1), M(5,2)"} \\
\\ 
\\ 
\texttt{¿Cuál es tu próximo comando?} \\
            }
        }
    }
    & \\ \\

    \theutterance \stepcounter{utterance}  
    & \multicolumn{4}{p{0.6\linewidth}}{
        \cellcolor[rgb]{0.9,0.9,0.9}{
            \makecell[{{p{\linewidth}}}]{
                \texttt{\tiny{[P1$\rangle$GM]}}
                \texttt{Vamos a analizar las diferencias entre ambos tableros. Mis posiciones y las suyas:} \\
\\ 
\texttt{Mis posiciones:} \\
\texttt{{-} P(2,5)} \\
\texttt{{-} O(3,4)} \\
\texttt{{-} T(2,1)} \\
\texttt{{-} S(2,6)} \\
\texttt{{-} D(7,2)} \\
\texttt{{-} A(1,2)} \\
\texttt{{-} M(3,6)} \\
\\ 
\texttt{Sus posiciones:} \\
\texttt{{-} P(3,1)} \\
\texttt{{-} O(7,1)} \\
\texttt{{-} T(7,6)} \\
\texttt{{-} S(1,2)} \\
\texttt{{-} D(5,3)} \\
\texttt{{-} A(5,1)} \\
\texttt{{-} M(5,2)} \\
\\ 
\texttt{Necesitamos acordar posiciones objetivo comunes. Como ambos tenemos que mover objetos, propongo un conjunto de coordenadas que podamos alcanzar. Voy a sugerirle una configuración que nos sirva a ambos.} \\
\\ 
\texttt{DECIR: Propongo las siguientes posiciones objetivo: P(3,3), O(4,3), T(3,1), S(3,5), D(5,3), A(4,1), M(4,5). ¿Qué te parece? ¿O prefieres sugerir tú otras coordenadas?} \\
            }
        }
    }
    & & \\ \\

    \theutterance \stepcounter{utterance}  
    & & & \multicolumn{2}{p{0.3\linewidth}}{
        \cellcolor[rgb]{0.9,0.9,0.9}{
            \makecell[{{p{\linewidth}}}]{
                \texttt{\tiny{[GM$|$GM]}}
                \texttt{Vamos a analizar las diferencias entre ambos tableros. Mis posiciones y las suyas:} \\
\\ 
\texttt{Mis posiciones:} \\
\texttt{{-} P(2,5)} \\
\texttt{{-} O(3,4)} \\
\texttt{{-} T(2,1)} \\
\texttt{{-} S(2,6)} \\
\texttt{{-} D(7,2)} \\
\texttt{{-} A(1,2)} \\
\texttt{{-} M(3,6)} \\
\\ 
\texttt{Sus posiciones:} \\
\texttt{{-} P(3,1)} \\
\texttt{{-} O(7,1)} \\
\texttt{{-} T(7,6)} \\
\texttt{{-} S(1,2)} \\
\texttt{{-} D(5,3)} \\
\texttt{{-} A(5,1)} \\
\texttt{{-} M(5,2)} \\
\\ 
\texttt{Necesitamos acordar posiciones objetivo comunes. Como ambos tenemos que mover objetos, propongo un conjunto de coordenadas que podamos alcanzar. Voy a sugerirle una configuración que nos sirva a ambos.} \\
\\ 
\texttt{DECIR: Propongo las siguientes posiciones objetivo: P(3,3), O(4,3), T(3,1), S(3,5), D(5,3), A(4,1), M(4,5). ¿Qué te parece? ¿O prefieres sugerir tú otras coordenadas?} \\
            }
        }
    }
    & & \\ \\

    \theutterance \stepcounter{utterance}  
    & & & \multicolumn{2}{p{0.3\linewidth}}{
        \cellcolor[rgb]{0.9,0.9,0.9}{
            \makecell[{{p{\linewidth}}}]{
                \texttt{\tiny{[GM$|$GM]}}
                \texttt{Invalid format: head is not empty: 'Vamos a analizar las diferencias entre ambos tableros. Mis posiciones y las suyas:} \\
\\ 
\texttt{Mis posiciones:} \\
\texttt{{-} P(2,5)} \\
\texttt{{-} O(3,4)} \\
\texttt{{-} T(2,1)} \\
\texttt{{-} S(2,6)} \\
\texttt{{-} D(7,2)} \\
\texttt{{-} A(1,2)} \\
\texttt{{-} M(3,6)} \\
\\ 
\texttt{Sus posiciones:} \\
\texttt{{-} P(3,1)} \\
\texttt{{-} O(7,1)} \\
\texttt{{-} T(7,6)} \\
\texttt{{-} S(1,2)} \\
\texttt{{-} D(5,3)} \\
\texttt{{-} A(5,1)} \\
\texttt{{-} M(5,2)} \\
\\ 
\texttt{Necesitamos acordar posiciones objetivo comunes. Como ambos tenemos que mover objetos, propongo un conjunto de coordenadas que podamos alcanzar. Voy a sugerirle una configuración que nos sirva a ambos.} \\
\\ 
\texttt{'} \\
            }
        }
    }
    & & \\ \\

    \theutterance \stepcounter{utterance}  
    & & \multicolumn{4}{p{0.6\linewidth}}{
        \cellcolor[rgb]{0.9,0.9,0.9}{
            \makecell[{{p{\linewidth}}}]{
                \texttt{\tiny{[P1$\langle$GM]}}
                \texttt{Penalización: ¡Tu mensaje no debe contener nada antes del comando!} \\
\\ 
\texttt{Asegúrate de que tu respuesta solo contenga `DECIR: <MENSAJE>` o `MOVERR: <OBJETO>, (<X>, <Y>)`, ¡y nada más!} \\
\\ 
\texttt{Han acumulado colectivamente 3 de 16 penalizaciones.} \\
\texttt{¡Por favor, inténtalo de nuevo!} \\
            }
        }
    }
    & \\ \\

    \theutterance \stepcounter{utterance}  
    & \multicolumn{4}{p{0.6\linewidth}}{
        \cellcolor[rgb]{0.9,0.9,0.9}{
            \makecell[{{p{\linewidth}}}]{
                \texttt{\tiny{[P1$\rangle$GM]}}
                \texttt{DECIR: Propongo estas coordenadas objetivo: P(3,3), O(4,3), T(3,1), S(3,5), D(5,3), A(4,1), M(4,5). ¿Te parece bien o sugieres otras?} \\
            }
        }
    }
    & & \\ \\

    \theutterance \stepcounter{utterance}  
    & & & \multicolumn{2}{p{0.3\linewidth}}{
        \cellcolor[rgb]{0.9,0.9,0.9}{
            \makecell[{{p{\linewidth}}}]{
                \texttt{\tiny{[GM$|$GM]}}
                \texttt{DECIR: Propongo estas coordenadas objetivo: P(3,3), O(4,3), T(3,1), S(3,5), D(5,3), A(4,1), M(4,5). ¿Te parece bien o sugieres otras?} \\
            }
        }
    }
    & & \\ \\

    \theutterance \stepcounter{utterance}  
    & & \multicolumn{4}{p{0.6\linewidth}}{
        \cellcolor[rgb]{0.9,0.9,0.9}{
            \makecell[{{p{\linewidth}}}]{
                \texttt{\tiny{[GM$\rangle$P2]}}
                \texttt{Tu mensaje ha sido transmitido al otro jugador.} \\
\\ 
\\ 
\texttt{Actualmente estás jugando la ronda 2 de un máximo de 28.} \\
\texttt{Han acumulado colectivamente 3 de 16 penalizaciones.} \\
\texttt{El otro jugador envió este mensaje:} \\
\\ 
\\ 
\texttt{"Propongo estas coordenadas objetivo: P(3,3), O(4,3), T(3,1), S(3,5), D(5,3), A(4,1), M(4,5). ¿Te parece bien o sugieres otras?"} \\
\\ 
\\ 
\texttt{¿Cuál es tu próximo comando?} \\
            }
        }
    }
    & \\ \\

    \theutterance \stepcounter{utterance}  
    & & & \multicolumn{4}{p{0.6\linewidth}}{
        \cellcolor[rgb]{0.9,0.9,0.9}{
            \makecell[{{p{\linewidth}}}]{
                \texttt{\tiny{[GM$\langle$P2]}}
                \texttt{Veamos si nuestras posiciones actuales coinciden con las propuestas o si hay conflictos.} \\
\\ 
\texttt{Mis posiciones actuales:} \\
\texttt{{-} P(3,1)} \\
\texttt{{-} O(7,1)} \\
\texttt{{-} T(7,6)} \\
\texttt{{-} S(1,2)} \\
\texttt{{-} D(5,3)} \\
\texttt{{-} A(5,1)} \\
\texttt{{-} M(5,2)} \\
\\ 
\texttt{Posiciones objetivo propuestas:} \\
\texttt{{-} P(3,3)} \\
\texttt{{-} O(4,3)} \\
\texttt{{-} T(3,1)} \\
\texttt{{-} S(3,5)} \\
\texttt{{-} D(5,3)} \\
\texttt{{-} A(4,1)} \\
\texttt{{-} M(4,5)} \\
\\ 
\texttt{Tengo D en (5,3) que ya coincide con el objetivo, lo cual es bueno.} \\
\\ 
\texttt{Revisando posibles conflictos: T(3,1) está vacío en mi tablero. A(4,1) está vacío. P(3,3) está vacío. O(4,3) está vacío. S(3,5) está vacío. M(4,5) está vacío. Parece que las posiciones propuestas están libres en mi tablero.} \\
\\ 
\texttt{La propuesta parece viable. Acepto.} \\
\\ 
\texttt{```} \\
\texttt{DECIR: Me parece bien, acepto las coordenadas objetivo: P(3,3), O(4,3), T(3,1), S(3,5), D(5,3), A(4,1), M(4,5). Empecemos a mover.} \\
            }
        }
    }
     \\ \\

    \theutterance \stepcounter{utterance}  
    & & & \multicolumn{2}{p{0.3\linewidth}}{
        \cellcolor[rgb]{0.9,0.9,0.9}{
            \makecell[{{p{\linewidth}}}]{
                \texttt{\tiny{[GM$|$GM]}}
                \texttt{Veamos si nuestras posiciones actuales coinciden con las propuestas o si hay conflictos.} \\
\\ 
\texttt{Mis posiciones actuales:} \\
\texttt{{-} P(3,1)} \\
\texttt{{-} O(7,1)} \\
\texttt{{-} T(7,6)} \\
\texttt{{-} S(1,2)} \\
\texttt{{-} D(5,3)} \\
\texttt{{-} A(5,1)} \\
\texttt{{-} M(5,2)} \\
\\ 
\texttt{Posiciones objetivo propuestas:} \\
\texttt{{-} P(3,3)} \\
\texttt{{-} O(4,3)} \\
\texttt{{-} T(3,1)} \\
\texttt{{-} S(3,5)} \\
\texttt{{-} D(5,3)} \\
\texttt{{-} A(4,1)} \\
\texttt{{-} M(4,5)} \\
\\ 
\texttt{Tengo D en (5,3) que ya coincide con el objetivo, lo cual es bueno.} \\
\\ 
\texttt{Revisando posibles conflictos: T(3,1) está vacío en mi tablero. A(4,1) está vacío. P(3,3) está vacío. O(4,3) está vacío. S(3,5) está vacío. M(4,5) está vacío. Parece que las posiciones propuestas están libres en mi tablero.} \\
\\ 
\texttt{La propuesta parece viable. Acepto.} \\
\\ 
\texttt{```} \\
\texttt{DECIR: Me parece bien, acepto las coordenadas objetivo: P(3,3), O(4,3), T(3,1), S(3,5), D(5,3), A(4,1), M(4,5). Empecemos a mover.} \\
            }
        }
    }
    & & \\ \\

    \theutterance \stepcounter{utterance}  
    & & & \multicolumn{2}{p{0.3\linewidth}}{
        \cellcolor[rgb]{0.9,0.9,0.9}{
            \makecell[{{p{\linewidth}}}]{
                \texttt{\tiny{[GM$|$GM]}}
                \texttt{Invalid format: head is not empty: 'Veamos si nuestras posiciones actuales coinciden con las propuestas o si hay conflictos.} \\
\\ 
\texttt{Mis posiciones actuales:} \\
\texttt{{-} P(3,1)} \\
\texttt{{-} O(7,1)} \\
\texttt{{-} T(7,6)} \\
\texttt{{-} S(1,2)} \\
\texttt{{-} D(5,3)} \\
\texttt{{-} A(5,1)} \\
\texttt{{-} M(5,2)} \\
\\ 
\texttt{Posiciones objetivo propuestas:} \\
\texttt{{-} P(3,3)} \\
\texttt{{-} O(4,3)} \\
\texttt{{-} T(3,1)} \\
\texttt{{-} S(3,5)} \\
\texttt{{-} D(5,3)} \\
\texttt{{-} A(4,1)} \\
\texttt{{-} M(4,5)} \\
\\ 
\texttt{Tengo D en (5,3) que ya coincide con el objetivo, lo cual es bueno.} \\
\\ 
\texttt{Revisando posibles conflictos: T(3,1) está vacío en mi tablero. A(4,1) está vacío. P(3,3) está vacío. O(4,3) está vacío. S(3,5) está vacío. M(4,5) está vacío. Parece que las posiciones propuestas están libres en mi tablero.} \\
\\ 
\texttt{La propuesta parece viable. Acepto.} \\
\\ 
\\ 
\texttt{'} \\
            }
        }
    }
    & & \\ \\

    \theutterance \stepcounter{utterance}  
    & & \multicolumn{4}{p{0.6\linewidth}}{
        \cellcolor[rgb]{0.9,0.9,0.9}{
            \makecell[{{p{\linewidth}}}]{
                \texttt{\tiny{[GM$\rangle$P2]}}
                \texttt{Penalización: ¡Tu mensaje no debe contener nada antes del comando!} \\
\\ 
\texttt{Asegúrate de que tu respuesta solo contenga `DECIR: <MENSAJE>` o `MOVERR: <OBJETO>, (<X>, <Y>)`, ¡y nada más!} \\
\\ 
\texttt{Han acumulado colectivamente 4 de 16 penalizaciones.} \\
\texttt{¡Por favor, inténtalo de nuevo!} \\
            }
        }
    }
    & \\ \\

    \theutterance \stepcounter{utterance}  
    & & & \multicolumn{4}{p{0.6\linewidth}}{
        \cellcolor[rgb]{0.9,0.9,0.9}{
            \makecell[{{p{\linewidth}}}]{
                \texttt{\tiny{[GM$\langle$P2]}}
                \texttt{DECIR: Me parece bien, acepto las coordenadas objetivo: P(3,3), O(4,3), T(3,1), S(3,5), D(5,3), A(4,1), M(4,5). Empecemos a mover.} \\
            }
        }
    }
     \\ \\

    \theutterance \stepcounter{utterance}  
    & & & \multicolumn{2}{p{0.3\linewidth}}{
        \cellcolor[rgb]{0.9,0.9,0.9}{
            \makecell[{{p{\linewidth}}}]{
                \texttt{\tiny{[GM$|$GM]}}
                \texttt{DECIR: Me parece bien, acepto las coordenadas objetivo: P(3,3), O(4,3), T(3,1), S(3,5), D(5,3), A(4,1), M(4,5). Empecemos a mover.} \\
            }
        }
    }
    & & \\ \\

    \theutterance \stepcounter{utterance}  
    & & \multicolumn{4}{p{0.6\linewidth}}{
        \cellcolor[rgb]{0.9,0.9,0.9}{
            \makecell[{{p{\linewidth}}}]{
                \texttt{\tiny{[P1$\langle$GM]}}
                \texttt{Tu mensaje ha sido transmitido al otro jugador.} \\
\\ 
\\ 
\texttt{Actualmente estás jugando la ronda 2 de un máximo de 28.} \\
\texttt{Han acumulado colectivamente 4 de 16 penalizaciones.} \\
\texttt{El otro jugador envió este mensaje:} \\
\\ 
\\ 
\texttt{"Me parece bien, acepto las coordenadas objetivo: P(3,3), O(4,3), T(3,1), S(3,5), D(5,3), A(4,1), M(4,5). Empecemos a mover."} \\
\\ 
\\ 
\texttt{¿Cuál es tu próximo comando?} \\
            }
        }
    }
    & \\ \\

    \theutterance \stepcounter{utterance}  
    & \multicolumn{4}{p{0.6\linewidth}}{
        \cellcolor[rgb]{0.9,0.9,0.9}{
            \makecell[{{p{\linewidth}}}]{
                \texttt{\tiny{[P1$\rangle$GM]}}
                \texttt{Primero verifico si alguno de mis objetos ya está en la posición objetivo acordada:} \\
\\ 
\texttt{Objetivos acordados:} \\
\texttt{{-} P(3,3) {-} Mi P está en (2,5), necesito moverlo} \\
\texttt{{-} O(4,3) {-} Mi O está en (3,4), necesito moverlo} \\
\texttt{{-} T(3,1) {-} Mi T está en (2,1), necesito moverlo} \\
\texttt{{-} S(3,5) {-} Mi S está en (2,6), necesito moverlo} \\
\texttt{{-} D(5,3) {-} Mi D está en (7,2), necesito moverlo} \\
\texttt{{-} A(4,1) {-} Mi A está en (1,2), necesito moverlo} \\
\texttt{{-} M(4,5) {-} Mi M está en (3,6), necesito moverlo} \\
\\ 
\texttt{Ninguno está en la posición correcta aún. Voy a empezar moviendo los objetos. Primero verificaré que las celdas destino estén vacías en mi tablero. Revisando mi cuadrícula inicial:} \\
\\ 
\texttt{Celda (3,1): tiene T actualmente, pero T está en (2,1). La celda (3,1) parece vacía. Moveré T a (3,1).} \\
\\ 
\texttt{MOVERR: T, (3,1)} \\
            }
        }
    }
    & & \\ \\

    \theutterance \stepcounter{utterance}  
    & & & \multicolumn{2}{p{0.3\linewidth}}{
        \cellcolor[rgb]{0.9,0.9,0.9}{
            \makecell[{{p{\linewidth}}}]{
                \texttt{\tiny{[GM$|$GM]}}
                \texttt{Primero verifico si alguno de mis objetos ya está en la posición objetivo acordada:} \\
\\ 
\texttt{Objetivos acordados:} \\
\texttt{{-} P(3,3) {-} Mi P está en (2,5), necesito moverlo} \\
\texttt{{-} O(4,3) {-} Mi O está en (3,4), necesito moverlo} \\
\texttt{{-} T(3,1) {-} Mi T está en (2,1), necesito moverlo} \\
\texttt{{-} S(3,5) {-} Mi S está en (2,6), necesito moverlo} \\
\texttt{{-} D(5,3) {-} Mi D está en (7,2), necesito moverlo} \\
\texttt{{-} A(4,1) {-} Mi A está en (1,2), necesito moverlo} \\
\texttt{{-} M(4,5) {-} Mi M está en (3,6), necesito moverlo} \\
\\ 
\texttt{Ninguno está en la posición correcta aún. Voy a empezar moviendo los objetos. Primero verificaré que las celdas destino estén vacías en mi tablero. Revisando mi cuadrícula inicial:} \\
\\ 
\texttt{Celda (3,1): tiene T actualmente, pero T está en (2,1). La celda (3,1) parece vacía. Moveré T a (3,1).} \\
\\ 
\texttt{MOVERR: T, (3,1)} \\
            }
        }
    }
    & & \\ \\

    \theutterance \stepcounter{utterance}  
    & & & \multicolumn{2}{p{0.3\linewidth}}{
        \cellcolor[rgb]{0.9,0.9,0.9}{
            \makecell[{{p{\linewidth}}}]{
                \texttt{\tiny{[GM$|$GM]}}
                \texttt{Invalid response format} \\
            }
        }
    }
    & & \\ \\

    \theutterance \stepcounter{utterance}  
    & & \multicolumn{4}{p{0.6\linewidth}}{
        \cellcolor[rgb]{0.9,0.9,0.9}{
            \makecell[{{p{\linewidth}}}]{
                \texttt{\tiny{[P1$\langle$GM]}}
                \texttt{Penalización: ¡Tu mensaje no tiene el formato esperado!} \\
\\ 
\texttt{Asegúrate de que tu respuesta solo contenga `DECIR: <MENSAJE>` o `MOVERR: <OBJETO>, (<X>, <Y>)`, ¡y nada más!} \\
\\ 
\texttt{Han acumulado colectivamente 5 de 16 penalizaciones.} \\
\texttt{¡Por favor, inténtalo de nuevo!} \\
            }
        }
    }
    & \\ \\

    \theutterance \stepcounter{utterance}  
    & \multicolumn{4}{p{0.6\linewidth}}{
        \cellcolor[rgb]{0.9,0.9,0.9}{
            \makecell[{{p{\linewidth}}}]{
                \texttt{\tiny{[P1$\rangle$GM]}}
                \texttt{MOVERR: T, (3,1)} \\
            }
        }
    }
    & & \\ \\

    \theutterance \stepcounter{utterance}  
    & & & \multicolumn{2}{p{0.3\linewidth}}{
        \cellcolor[rgb]{0.9,0.9,0.9}{
            \makecell[{{p{\linewidth}}}]{
                \texttt{\tiny{[GM$|$GM]}}
                \texttt{MOVERR: T, (3,1)} \\
            }
        }
    }
    & & \\ \\

    \theutterance \stepcounter{utterance}  
    & & & \multicolumn{2}{p{0.3\linewidth}}{
        \cellcolor[rgb]{0.9,0.9,0.9}{
            \makecell[{{p{\linewidth}}}]{
                \texttt{\tiny{[GM$|$GM]}}
                \texttt{Invalid response format} \\
            }
        }
    }
    & & \\ \\

    \theutterance \stepcounter{utterance}  
    & & \multicolumn{4}{p{0.6\linewidth}}{
        \cellcolor[rgb]{0.9,0.9,0.9}{
            \makecell[{{p{\linewidth}}}]{
                \texttt{\tiny{[P1$\langle$GM]}}
                \texttt{Penalización: ¡Tu mensaje no tiene el formato esperado!} \\
\\ 
\texttt{Asegúrate de que tu respuesta solo contenga `DECIR: <MENSAJE>` o `MOVERR: <OBJETO>, (<X>, <Y>)`, ¡y nada más!} \\
\\ 
\texttt{Han acumulado colectivamente 6 de 16 penalizaciones.} \\
\texttt{¡Por favor, inténtalo de nuevo!} \\
            }
        }
    }
    & \\ \\

    \theutterance \stepcounter{utterance}  
    & \multicolumn{4}{p{0.6\linewidth}}{
        \cellcolor[rgb]{0.9,0.9,0.9}{
            \makecell[{{p{\linewidth}}}]{
                \texttt{\tiny{[P1$\rangle$GM]}}
                \texttt{Veo que hay espacios después de la coma. Corrijo el formato.} \\
\\ 
\texttt{MOVERR: T,(3,1)} \\
            }
        }
    }
    & & \\ \\

    \theutterance \stepcounter{utterance}  
    & & & \multicolumn{2}{p{0.3\linewidth}}{
        \cellcolor[rgb]{0.9,0.9,0.9}{
            \makecell[{{p{\linewidth}}}]{
                \texttt{\tiny{[GM$|$GM]}}
                \texttt{Veo que hay espacios después de la coma. Corrijo el formato.} \\
\\ 
\texttt{MOVERR: T,(3,1)} \\
            }
        }
    }
    & & \\ \\

    \theutterance \stepcounter{utterance}  
    & & & \multicolumn{2}{p{0.3\linewidth}}{
        \cellcolor[rgb]{0.9,0.9,0.9}{
            \makecell[{{p{\linewidth}}}]{
                \texttt{\tiny{[GM$|$GM]}}
                \texttt{Invalid response format} \\
            }
        }
    }
    & & \\ \\

    \theutterance \stepcounter{utterance}  
    & & \multicolumn{4}{p{0.6\linewidth}}{
        \cellcolor[rgb]{0.9,0.9,0.9}{
            \makecell[{{p{\linewidth}}}]{
                \texttt{\tiny{[P1$\langle$GM]}}
                \texttt{Penalización: ¡Tu mensaje no tiene el formato esperado!} \\
\\ 
\texttt{Asegúrate de que tu respuesta solo contenga `DECIR: <MENSAJE>` o `MOVERR: <OBJETO>, (<X>, <Y>)`, ¡y nada más!} \\
\\ 
\texttt{Han acumulado colectivamente 7 de 16 penalizaciones.} \\
\texttt{¡Por favor, inténtalo de nuevo!} \\
            }
        }
    }
    & \\ \\

    \theutterance \stepcounter{utterance}  
    & \multicolumn{4}{p{0.6\linewidth}}{
        \cellcolor[rgb]{0.9,0.9,0.9}{
            \makecell[{{p{\linewidth}}}]{
                \texttt{\tiny{[P1$\rangle$GM]}}
                \texttt{Revisando el formato correcto: `MOVERR: <OBJETO>, (<X>, <Y>)` {-} veo que requiere un espacio después de la coma y los paréntesis incluidos.} \\
\\ 
\texttt{MOVERR: T, (3,1)} \\
            }
        }
    }
    & & \\ \\

    \theutterance \stepcounter{utterance}  
    & & & \multicolumn{2}{p{0.3\linewidth}}{
        \cellcolor[rgb]{0.9,0.9,0.9}{
            \makecell[{{p{\linewidth}}}]{
                \texttt{\tiny{[GM$|$GM]}}
                \texttt{Revisando el formato correcto: `MOVERR: <OBJETO>, (<X>, <Y>)` {-} veo que requiere un espacio después de la coma y los paréntesis incluidos.} \\
\\ 
\texttt{MOVERR: T, (3,1)} \\
            }
        }
    }
    & & \\ \\

    \theutterance \stepcounter{utterance}  
    & & & \multicolumn{2}{p{0.3\linewidth}}{
        \cellcolor[rgb]{0.9,0.9,0.9}{
            \makecell[{{p{\linewidth}}}]{
                \texttt{\tiny{[GM$|$GM]}}
                \texttt{Invalid response format} \\
            }
        }
    }
    & & \\ \\

    \theutterance \stepcounter{utterance}  
    & & \multicolumn{4}{p{0.6\linewidth}}{
        \cellcolor[rgb]{0.9,0.9,0.9}{
            \makecell[{{p{\linewidth}}}]{
                \texttt{\tiny{[P1$\langle$GM]}}
                \texttt{Penalización: ¡Tu mensaje no tiene el formato esperado!} \\
\\ 
\texttt{Asegúrate de que tu respuesta solo contenga `DECIR: <MENSAJE>` o `MOVERR: <OBJETO>, (<X>, <Y>)`, ¡y nada más!} \\
\\ 
\texttt{Han acumulado colectivamente 8 de 16 penalizaciones.} \\
\texttt{¡Por favor, inténtalo de nuevo!} \\
            }
        }
    }
    & \\ \\

    \theutterance \stepcounter{utterance}  
    & \multicolumn{4}{p{0.6\linewidth}}{
        \cellcolor[rgb]{0.9,0.9,0.9}{
            \makecell[{{p{\linewidth}}}]{
                \texttt{\tiny{[P1$\rangle$GM]}}
                \texttt{Déjame revisar el formato exacto del ejemplo: `MOVERR: <OBJETO>, (<X>, <Y>)`. El problema puede ser que estoy incluyendo texto adicional sin darme cuenta. Lo intentaré con el formato más limpio posible.} \\
\\ 
\texttt{MOVERR: T,(3,1)} \\
            }
        }
    }
    & & \\ \\

    \theutterance \stepcounter{utterance}  
    & & & \multicolumn{2}{p{0.3\linewidth}}{
        \cellcolor[rgb]{0.9,0.9,0.9}{
            \makecell[{{p{\linewidth}}}]{
                \texttt{\tiny{[GM$|$GM]}}
                \texttt{Déjame revisar el formato exacto del ejemplo: `MOVERR: <OBJETO>, (<X>, <Y>)`. El problema puede ser que estoy incluyendo texto adicional sin darme cuenta. Lo intentaré con el formato más limpio posible.} \\
\\ 
\texttt{MOVERR: T,(3,1)} \\
            }
        }
    }
    & & \\ \\

    \theutterance \stepcounter{utterance}  
    & & & \multicolumn{2}{p{0.3\linewidth}}{
        \cellcolor[rgb]{0.9,0.9,0.9}{
            \makecell[{{p{\linewidth}}}]{
                \texttt{\tiny{[GM$|$GM]}}
                \texttt{Invalid response format} \\
            }
        }
    }
    & & \\ \\

    \theutterance \stepcounter{utterance}  
    & & \multicolumn{4}{p{0.6\linewidth}}{
        \cellcolor[rgb]{0.9,0.9,0.9}{
            \makecell[{{p{\linewidth}}}]{
                \texttt{\tiny{[P1$\langle$GM]}}
                \texttt{Penalización: ¡Tu mensaje no tiene el formato esperado!} \\
\\ 
\texttt{Asegúrate de que tu respuesta solo contenga `DECIR: <MENSAJE>` o `MOVERR: <OBJETO>, (<X>, <Y>)`, ¡y nada más!} \\
\\ 
\texttt{Han acumulado colectivamente 9 de 16 penalizaciones.} \\
\texttt{¡Por favor, inténtalo de nuevo!} \\
            }
        }
    }
    & \\ \\

    \theutterance \stepcounter{utterance}  
    & \multicolumn{4}{p{0.6\linewidth}}{
        \cellcolor[rgb]{0.9,0.9,0.9}{
            \makecell[{{p{\linewidth}}}]{
                \texttt{\tiny{[P1$\rangle$GM]}}
                \texttt{Quizás el formato correcto requiere un espacio exactamente como está en el ejemplo: después de MOVERR: hay un espacio, luego el objeto, luego una coma, luego un espacio, luego el paréntesis. Voy a intentar copiando el formato del ejemplo original donde decía "MOVERR: <OBJETO>, (<X>, <Y>)" interpretándolo literalmente con el espacio después de la coma.} \\
\\ 
\texttt{MOVERR: T,(3,1)} \\
            }
        }
    }
    & & \\ \\

    \theutterance \stepcounter{utterance}  
    & & & \multicolumn{2}{p{0.3\linewidth}}{
        \cellcolor[rgb]{0.9,0.9,0.9}{
            \makecell[{{p{\linewidth}}}]{
                \texttt{\tiny{[GM$|$GM]}}
                \texttt{Quizás el formato correcto requiere un espacio exactamente como está en el ejemplo: después de MOVERR: hay un espacio, luego el objeto, luego una coma, luego un espacio, luego el paréntesis. Voy a intentar copiando el formato del ejemplo original donde decía "MOVERR: <OBJETO>, (<X>, <Y>)" interpretándolo literalmente con el espacio después de la coma.} \\
\\ 
\texttt{MOVERR: T,(3,1)} \\
            }
        }
    }
    & & \\ \\

    \theutterance \stepcounter{utterance}  
    & & & \multicolumn{2}{p{0.3\linewidth}}{
        \cellcolor[rgb]{0.9,0.9,0.9}{
            \makecell[{{p{\linewidth}}}]{
                \texttt{\tiny{[GM$|$GM]}}
                \texttt{Invalid response format} \\
            }
        }
    }
    & & \\ \\

    \theutterance \stepcounter{utterance}  
    & & \multicolumn{4}{p{0.6\linewidth}}{
        \cellcolor[rgb]{0.9,0.9,0.9}{
            \makecell[{{p{\linewidth}}}]{
                \texttt{\tiny{[P1$\langle$GM]}}
                \texttt{Penalización: ¡Tu mensaje no tiene el formato esperado!} \\
\\ 
\texttt{Asegúrate de que tu respuesta solo contenga `DECIR: <MENSAJE>` o `MOVERR: <OBJETO>, (<X>, <Y>)`, ¡y nada más!} \\
\\ 
\texttt{Han acumulado colectivamente 10 de 16 penalizaciones.} \\
\texttt{¡Por favor, inténtalo de nuevo!} \\
            }
        }
    }
    & \\ \\

    \theutterance \stepcounter{utterance}  
    & \multicolumn{4}{p{0.6\linewidth}}{
        \cellcolor[rgb]{0.9,0.9,0.9}{
            \makecell[{{p{\linewidth}}}]{
                \texttt{\tiny{[P1$\rangle$GM]}}
                \texttt{Voy a intentar con el formato exactamente como aparece en las instrucciones, con el objeto entre < >. Quizás esa es la sintaxis correcta.} \\
\\ 
\texttt{MOVERR: <T>, (<3>, <1>)} \\
            }
        }
    }
    & & \\ \\

    \theutterance \stepcounter{utterance}  
    & & & \multicolumn{2}{p{0.3\linewidth}}{
        \cellcolor[rgb]{0.9,0.9,0.9}{
            \makecell[{{p{\linewidth}}}]{
                \texttt{\tiny{[GM$|$GM]}}
                \texttt{Voy a intentar con el formato exactamente como aparece en las instrucciones, con el objeto entre < >. Quizás esa es la sintaxis correcta.} \\
\\ 
\texttt{MOVERR: <T>, (<3>, <1>)} \\
            }
        }
    }
    & & \\ \\

    \theutterance \stepcounter{utterance}  
    & & & \multicolumn{2}{p{0.3\linewidth}}{
        \cellcolor[rgb]{0.9,0.9,0.9}{
            \makecell[{{p{\linewidth}}}]{
                \texttt{\tiny{[GM$|$GM]}}
                \texttt{Invalid response format} \\
            }
        }
    }
    & & \\ \\

    \theutterance \stepcounter{utterance}  
    & & \multicolumn{4}{p{0.6\linewidth}}{
        \cellcolor[rgb]{0.9,0.9,0.9}{
            \makecell[{{p{\linewidth}}}]{
                \texttt{\tiny{[P1$\langle$GM]}}
                \texttt{Penalización: ¡Tu mensaje no tiene el formato esperado!} \\
\\ 
\texttt{Asegúrate de que tu respuesta solo contenga `DECIR: <MENSAJE>` o `MOVERR: <OBJETO>, (<X>, <Y>)`, ¡y nada más!} \\
\\ 
\texttt{Han acumulado colectivamente 11 de 16 penalizaciones.} \\
\texttt{¡Por favor, inténtalo de nuevo!} \\
            }
        }
    }
    & \\ \\

    \theutterance \stepcounter{utterance}  
    & \multicolumn{4}{p{0.6\linewidth}}{
        \cellcolor[rgb]{0.9,0.9,0.9}{
            \makecell[{{p{\linewidth}}}]{
                \texttt{\tiny{[P1$\rangle$GM]}}
                \texttt{Intentaré simplemente con el texto más mínimo posible y el formato exacto.} \\
\\ 
\texttt{MOVERR: T, (3, 1)} \\
            }
        }
    }
    & & \\ \\

    \theutterance \stepcounter{utterance}  
    & & & \multicolumn{2}{p{0.3\linewidth}}{
        \cellcolor[rgb]{0.9,0.9,0.9}{
            \makecell[{{p{\linewidth}}}]{
                \texttt{\tiny{[GM$|$GM]}}
                \texttt{Intentaré simplemente con el texto más mínimo posible y el formato exacto.} \\
\\ 
\texttt{MOVERR: T, (3, 1)} \\
            }
        }
    }
    & & \\ \\

    \theutterance \stepcounter{utterance}  
    & & & \multicolumn{2}{p{0.3\linewidth}}{
        \cellcolor[rgb]{0.9,0.9,0.9}{
            \makecell[{{p{\linewidth}}}]{
                \texttt{\tiny{[GM$|$GM]}}
                \texttt{Invalid response format} \\
            }
        }
    }
    & & \\ \\

    \theutterance \stepcounter{utterance}  
    & & \multicolumn{4}{p{0.6\linewidth}}{
        \cellcolor[rgb]{0.9,0.9,0.9}{
            \makecell[{{p{\linewidth}}}]{
                \texttt{\tiny{[P1$\langle$GM]}}
                \texttt{Penalización: ¡Tu mensaje no tiene el formato esperado!} \\
\\ 
\texttt{Asegúrate de que tu respuesta solo contenga `DECIR: <MENSAJE>` o `MOVERR: <OBJETO>, (<X>, <Y>)`, ¡y nada más!} \\
\\ 
\texttt{Han acumulado colectivamente 12 de 16 penalizaciones.} \\
\texttt{¡Por favor, inténtalo de nuevo!} \\
            }
        }
    }
    & \\ \\

    \theutterance \stepcounter{utterance}  
    & \multicolumn{4}{p{0.6\linewidth}}{
        \cellcolor[rgb]{0.9,0.9,0.9}{
            \makecell[{{p{\linewidth}}}]{
                \texttt{\tiny{[P1$\rangle$GM]}}
                \texttt{Voy a enviar un mensaje al otro jugador para preguntarle cómo está ejecutando sus movimientos correctamente.} \\
\\ 
\texttt{DECIR: Estoy teniendo problemas con el formato del comando MOVERR. ¿Podrías mostrarme exactamente cómo escribes tú un movimiento exitoso?} \\
            }
        }
    }
    & & \\ \\

    \theutterance \stepcounter{utterance}  
    & & & \multicolumn{2}{p{0.3\linewidth}}{
        \cellcolor[rgb]{0.9,0.9,0.9}{
            \makecell[{{p{\linewidth}}}]{
                \texttt{\tiny{[GM$|$GM]}}
                \texttt{Voy a enviar un mensaje al otro jugador para preguntarle cómo está ejecutando sus movimientos correctamente.} \\
\\ 
\texttt{DECIR: Estoy teniendo problemas con el formato del comando MOVERR. ¿Podrías mostrarme exactamente cómo escribes tú un movimiento exitoso?} \\
            }
        }
    }
    & & \\ \\

    \theutterance \stepcounter{utterance}  
    & & & \multicolumn{2}{p{0.3\linewidth}}{
        \cellcolor[rgb]{0.9,0.9,0.9}{
            \makecell[{{p{\linewidth}}}]{
                \texttt{\tiny{[GM$|$GM]}}
                \texttt{Invalid format: head is not empty: 'Voy a enviar un mensaje al otro jugador para preguntarle cómo está ejecutando sus movimientos correctamente.} \\
\\ 
\texttt{'} \\
            }
        }
    }
    & & \\ \\

    \theutterance \stepcounter{utterance}  
    & & \multicolumn{4}{p{0.6\linewidth}}{
        \cellcolor[rgb]{0.9,0.9,0.9}{
            \makecell[{{p{\linewidth}}}]{
                \texttt{\tiny{[P1$\langle$GM]}}
                \texttt{Penalización: ¡Tu mensaje no debe contener nada antes del comando!} \\
\\ 
\texttt{Asegúrate de que tu respuesta solo contenga `DECIR: <MENSAJE>` o `MOVERR: <OBJETO>, (<X>, <Y>)`, ¡y nada más!} \\
\\ 
\texttt{Han acumulado colectivamente 13 de 16 penalizaciones.} \\
\texttt{¡Por favor, inténtalo de nuevo!} \\
            }
        }
    }
    & \\ \\

    \theutterance \stepcounter{utterance}  
    & \multicolumn{4}{p{0.6\linewidth}}{
        \cellcolor[rgb]{0.9,0.9,0.9}{
            \makecell[{{p{\linewidth}}}]{
                \texttt{\tiny{[P1$\rangle$GM]}}
                \texttt{DECIR: No logro hacer funcionar MOVERR, ¿puedes mostrarme un ejemplo exacto de cómo escribes un movimiento exitoso?} \\
            }
        }
    }
    & & \\ \\

    \theutterance \stepcounter{utterance}  
    & & & \multicolumn{2}{p{0.3\linewidth}}{
        \cellcolor[rgb]{0.9,0.9,0.9}{
            \makecell[{{p{\linewidth}}}]{
                \texttt{\tiny{[GM$|$GM]}}
                \texttt{DECIR: No logro hacer funcionar MOVERR, ¿puedes mostrarme un ejemplo exacto de cómo escribes un movimiento exitoso?} \\
            }
        }
    }
    & & \\ \\

    \theutterance \stepcounter{utterance}  
    & & \multicolumn{4}{p{0.6\linewidth}}{
        \cellcolor[rgb]{0.9,0.9,0.9}{
            \makecell[{{p{\linewidth}}}]{
                \texttt{\tiny{[GM$\rangle$P2]}}
                \texttt{Tu mensaje ha sido transmitido al otro jugador.} \\
\\ 
\\ 
\texttt{Actualmente estás jugando la ronda 3 de un máximo de 28.} \\
\texttt{Han acumulado colectivamente 13 de 16 penalizaciones.} \\
\texttt{El otro jugador envió este mensaje:} \\
\\ 
\\ 
\texttt{"No logro hacer funcionar MOVERR, ¿puedes mostrarme un ejemplo exacto de cómo escribes un movimiento exitoso?"} \\
\\ 
\\ 
\texttt{¿Cuál es tu próximo comando?} \\
            }
        }
    }
    & \\ \\

    \theutterance \stepcounter{utterance}  
    & & & \multicolumn{4}{p{0.6\linewidth}}{
        \cellcolor[rgb]{0.9,0.9,0.9}{
            \makecell[{{p{\linewidth}}}]{
                \texttt{\tiny{[GM$\langle$P2]}}
                \texttt{Voy a mover el objeto D a (5,3), que ya está ahí, pero así muestro el formato. Mejor voy a mover un objeto que necesite mover. Por ejemplo, moveré A de (5,1) a (4,1).} \\
\\ 
\texttt{```} \\
\texttt{MOVERR: A, (4, 1)} \\
\texttt{```} \\
            }
        }
    }
     \\ \\

    \theutterance \stepcounter{utterance}  
    & & & \multicolumn{2}{p{0.3\linewidth}}{
        \cellcolor[rgb]{0.9,0.9,0.9}{
            \makecell[{{p{\linewidth}}}]{
                \texttt{\tiny{[GM$|$GM]}}
                \texttt{Voy a mover el objeto D a (5,3), que ya está ahí, pero así muestro el formato. Mejor voy a mover un objeto que necesite mover. Por ejemplo, moveré A de (5,1) a (4,1).} \\
\\ 
\texttt{```} \\
\texttt{MOVERR: A, (4, 1)} \\
\texttt{```} \\
            }
        }
    }
    & & \\ \\

    \theutterance \stepcounter{utterance}  
    & & & \multicolumn{2}{p{0.3\linewidth}}{
        \cellcolor[rgb]{0.9,0.9,0.9}{
            \makecell[{{p{\linewidth}}}]{
                \texttt{\tiny{[GM$|$GM]}}
                \texttt{Invalid response format} \\
            }
        }
    }
    & & \\ \\

    \theutterance \stepcounter{utterance}  
    & & \multicolumn{4}{p{0.6\linewidth}}{
        \cellcolor[rgb]{0.9,0.9,0.9}{
            \makecell[{{p{\linewidth}}}]{
                \texttt{\tiny{[GM$\rangle$P2]}}
                \texttt{Penalización: ¡Tu mensaje no tiene el formato esperado!} \\
\\ 
\texttt{Asegúrate de que tu respuesta solo contenga `DECIR: <MENSAJE>` o `MOVERR: <OBJETO>, (<X>, <Y>)`, ¡y nada más!} \\
\\ 
\texttt{Han acumulado colectivamente 14 de 16 penalizaciones.} \\
\texttt{¡Por favor, inténtalo de nuevo!} \\
            }
        }
    }
    & \\ \\

    \theutterance \stepcounter{utterance}  
    & & & \multicolumn{4}{p{0.6\linewidth}}{
        \cellcolor[rgb]{0.9,0.9,0.9}{
            \makecell[{{p{\linewidth}}}]{
                \texttt{\tiny{[GM$\langle$P2]}}
                \texttt{MOVERR: A, (4, 1)} \\
            }
        }
    }
     \\ \\

    \theutterance \stepcounter{utterance}  
    & & & \multicolumn{2}{p{0.3\linewidth}}{
        \cellcolor[rgb]{0.9,0.9,0.9}{
            \makecell[{{p{\linewidth}}}]{
                \texttt{\tiny{[GM$|$GM]}}
                \texttt{MOVERR: A, (4, 1)} \\
            }
        }
    }
    & & \\ \\

    \theutterance \stepcounter{utterance}  
    & & & \multicolumn{2}{p{0.3\linewidth}}{
        \cellcolor[rgb]{0.9,0.9,0.9}{
            \makecell[{{p{\linewidth}}}]{
                \texttt{\tiny{[GM$|$GM]}}
                \texttt{Invalid response format} \\
            }
        }
    }
    & & \\ \\

    \theutterance \stepcounter{utterance}  
    & & \multicolumn{4}{p{0.6\linewidth}}{
        \cellcolor[rgb]{0.9,0.9,0.9}{
            \makecell[{{p{\linewidth}}}]{
                \texttt{\tiny{[GM$\rangle$P2]}}
                \texttt{Penalización: ¡Tu mensaje no tiene el formato esperado!} \\
\\ 
\texttt{Asegúrate de que tu respuesta solo contenga `DECIR: <MENSAJE>` o `MOVERR: <OBJETO>, (<X>, <Y>)`, ¡y nada más!} \\
\\ 
\texttt{Han acumulado colectivamente 15 de 16 penalizaciones.} \\
\texttt{¡Por favor, inténtalo de nuevo!} \\
            }
        }
    }
    & \\ \\

    \theutterance \stepcounter{utterance}  
    & & & \multicolumn{4}{p{0.6\linewidth}}{
        \cellcolor[rgb]{0.9,0.9,0.9}{
            \makecell[{{p{\linewidth}}}]{
                \texttt{\tiny{[GM$\langle$P2]}}
                \texttt{DECIR: Ejemplo de movimiento: MOVERR: A, (4, 1)} \\
            }
        }
    }
     \\ \\

    \theutterance \stepcounter{utterance}  
    & & & \multicolumn{2}{p{0.3\linewidth}}{
        \cellcolor[rgb]{0.9,0.9,0.9}{
            \makecell[{{p{\linewidth}}}]{
                \texttt{\tiny{[GM$|$GM]}}
                \texttt{DECIR: Ejemplo de movimiento: MOVERR: A, (4, 1)} \\
            }
        }
    }
    & & \\ \\

    \theutterance \stepcounter{utterance}  
    & & \multicolumn{4}{p{0.6\linewidth}}{
        \cellcolor[rgb]{0.9,0.9,0.9}{
            \makecell[{{p{\linewidth}}}]{
                \texttt{\tiny{[P1$\langle$GM]}}
                \texttt{Tu mensaje ha sido transmitido al otro jugador.} \\
\\ 
\\ 
\texttt{Actualmente estás jugando la ronda 3 de un máximo de 28.} \\
\texttt{Han acumulado colectivamente 15 de 16 penalizaciones.} \\
\texttt{El otro jugador envió este mensaje:} \\
\\ 
\\ 
\texttt{"Ejemplo de movimiento: MOVERR: A, (4, 1)"} \\
\\ 
\\ 
\texttt{¿Cuál es tu próximo comando?} \\
            }
        }
    }
    & \\ \\

    \theutterance \stepcounter{utterance}  
    & \multicolumn{4}{p{0.6\linewidth}}{
        \cellcolor[rgb]{0.9,0.9,0.9}{
            \makecell[{{p{\linewidth}}}]{
                \texttt{\tiny{[P1$\rangle$GM]}}
                \texttt{¡Gracias! Ahora intentaré con el formato exacto que me compartió: `MOVERR: A, (4, 1)`} \\
\\ 
\texttt{MOVERR: T, (3, 1)} \\
            }
        }
    }
    & & \\ \\

    \theutterance \stepcounter{utterance}  
    & & & \multicolumn{2}{p{0.3\linewidth}}{
        \cellcolor[rgb]{0.9,0.9,0.9}{
            \makecell[{{p{\linewidth}}}]{
                \texttt{\tiny{[GM$|$GM]}}
                \texttt{¡Gracias! Ahora intentaré con el formato exacto que me compartió: `MOVERR: A, (4, 1)`} \\
\\ 
\texttt{MOVERR: T, (3, 1)} \\
            }
        }
    }
    & & \\ \\

    \theutterance \stepcounter{utterance}  
    & & & \multicolumn{2}{p{0.3\linewidth}}{
        \cellcolor[rgb]{0.9,0.9,0.9}{
            \makecell[{{p{\linewidth}}}]{
                \texttt{\tiny{[GM$|$GM]}}
                \texttt{Invalid response format} \\
            }
        }
    }
    & & \\ \\

    \theutterance \stepcounter{utterance}  
    & & \multicolumn{4}{p{0.6\linewidth}}{
        \cellcolor[rgb]{0.9,0.9,0.9}{
            \makecell[{{p{\linewidth}}}]{
                \texttt{\tiny{[P1$\langle$GM]}}
                \texttt{Penalización: ¡Tu mensaje no tiene el formato esperado!} \\
\\ 
\texttt{Asegúrate de que tu respuesta solo contenga `DECIR: <MENSAJE>` o `MOVERR: <OBJETO>, (<X>, <Y>)`, ¡y nada más!} \\
\\ 
\texttt{Han acumulado colectivamente 16 de 16 penalizaciones.} \\
\texttt{¡Por favor, inténtalo de nuevo!} \\
            }
        }
    }
    & \\ \\

    \theutterance \stepcounter{utterance}  
    & \multicolumn{4}{p{0.6\linewidth}}{
        \cellcolor[rgb]{0.9,0.9,0.9}{
            \makecell[{{p{\linewidth}}}]{
                \texttt{\tiny{[P1$\rangle$GM]}}
                \texttt{¡Hemos alcanzado las 16 penalizaciones! Según las reglas: "Si ambos jugadores acumulan más de 16 penalizaciones, ambos pierden el juego."} \\
\\ 
\texttt{Lamentablemente, hemos perdido el juego debido a las penalizaciones acumuladas. El formato de mis comandos no fue aceptado por el sistema a pesar de seguir el formato indicado. Pido disculpas por no haber logrado ejecutar los movimientos correctamente. ¿Quieres que intentemos algo diferente o podemos analizar qué salió mal?} \\
            }
        }
    }
    & & \\ \\

    \theutterance \stepcounter{utterance}  
    & & & \multicolumn{2}{p{0.3\linewidth}}{
        \cellcolor[rgb]{0.9,0.9,0.9}{
            \makecell[{{p{\linewidth}}}]{
                \texttt{\tiny{[GM$|$GM]}}
                \texttt{¡Hemos alcanzado las 16 penalizaciones! Según las reglas: "Si ambos jugadores acumulan más de 16 penalizaciones, ambos pierden el juego."} \\
\\ 
\texttt{Lamentablemente, hemos perdido el juego debido a las penalizaciones acumuladas. El formato de mis comandos no fue aceptado por el sistema a pesar de seguir el formato indicado. Pido disculpas por no haber logrado ejecutar los movimientos correctamente. ¿Quieres que intentemos algo diferente o podemos analizar qué salió mal?} \\
            }
        }
    }
    & & \\ \\

    \theutterance \stepcounter{utterance}  
    & & & \multicolumn{2}{p{0.3\linewidth}}{
        \cellcolor[rgb]{0.9,0.9,0.9}{
            \makecell[{{p{\linewidth}}}]{
                \texttt{\tiny{[GM$|$GM]}}
                \texttt{Invalid response format} \\
            }
        }
    }
    & & \\ \\

    \theutterance \stepcounter{utterance}  
    & & & \multicolumn{2}{p{0.3\linewidth}}{
        \cellcolor[rgb]{0.9,0.9,0.9}{
            \makecell[{{p{\linewidth}}}]{
                \texttt{\tiny{[GM$|$GM]}}
                \texttt{Maximum number of penalties exceeded} \\
            }
        }
    }
    & & \\ \\

    \theutterance \stepcounter{utterance}  
    & & & \multicolumn{2}{p{0.3\linewidth}}{
        \cellcolor[rgb]{0.9,0.9,0.9}{
            \makecell[{{p{\linewidth}}}]{
                \texttt{\tiny{[GM$|$GM]}}
                \texttt{Initial states:} \\
\texttt{```} \\
\texttt{Player 1:} \\
\texttt{ 1234567} \\
\texttt{╔═══════╗} \\
\texttt{║◌T◌◌◌◌◌║ 1} \\
\texttt{║A◌◌◌◌◌D║ 2} \\
\texttt{║◌◌◌◌◌┌─╢ 3} \\
\texttt{╟──┐O◌│◌║ 4} \\
\texttt{║◌P├──┼─╢ 5} \\
\texttt{║◌S│M◌│◌║ 6} \\
\texttt{║◌◌│◌◌│◌║ 7} \\
\texttt{╚══╧══╧═╝} \\
\\ 
\\ 
\texttt{Player 2:} \\
\texttt{ 1234567} \\
\texttt{╔═══════╗} \\
\texttt{║◌◌P◌A◌O║ 1} \\
\texttt{║S◌◌◌M◌◌║ 2} \\
\texttt{║◌◌◌◌D┌─╢ 3} \\
\texttt{╟──┐◌◌│◌║ 4} \\
\texttt{║◌◌├──┼─╢ 5} \\
\texttt{║◌◌│◌◌│T║ 6} \\
\texttt{║◌◌│◌◌│◌║ 7} \\
\texttt{╚══╧══╧═╝} \\
\\ 
\texttt{```} \\
            }
        }
    }
    & & \\ \\

    \theutterance \stepcounter{utterance}  
    & & & \multicolumn{2}{p{0.3\linewidth}}{
        \cellcolor[rgb]{0.9,0.9,0.9}{
            \makecell[{{p{\linewidth}}}]{
                \texttt{\tiny{[GM$|$GM]}}
                \texttt{End states:} \\
\texttt{```} \\
\texttt{Player 1:} \\
\texttt{ 1234567} \\
\texttt{╔═══════╗} \\
\texttt{║◌T◌◌◌◌◌║ 1} \\
\texttt{║A◌◌◌◌◌D║ 2} \\
\texttt{║◌◌◌◌◌┌─╢ 3} \\
\texttt{╟──┐O◌│◌║ 4} \\
\texttt{║◌P├──┼─╢ 5} \\
\texttt{║◌S│M◌│◌║ 6} \\
\texttt{║◌◌│◌◌│◌║ 7} \\
\texttt{╚══╧══╧═╝} \\
\\ 
\\ 
\texttt{Player 2:} \\
\texttt{ 1234567} \\
\texttt{╔═══════╗} \\
\texttt{║◌◌P◌A◌O║ 1} \\
\texttt{║S◌◌◌M◌◌║ 2} \\
\texttt{║◌◌◌◌D┌─╢ 3} \\
\texttt{╟──┐◌◌│◌║ 4} \\
\texttt{║◌◌├──┼─╢ 5} \\
\texttt{║◌◌│◌◌│T║ 6} \\
\texttt{║◌◌│◌◌│◌║ 7} \\
\texttt{╚══╧══╧═╝} \\
\\ 
\texttt{```} \\
            }
        }
    }
    & & \\ \\

    \theutterance \stepcounter{utterance}  
    & & & \multicolumn{2}{p{0.3\linewidth}}{
        \cellcolor[rgb]{0.9,0.9,0.9}{
            \makecell[{{p{\linewidth}}}]{
                \texttt{\tiny{[GM$|$GM]}}
                \texttt{* success: False} \\
\texttt{* lose: False} \\
\texttt{* aborted: True} \\
\texttt{{-}{-}{-}{-}{-}{-}{-}} \\
\texttt{* Shifts: 0.00} \\
\texttt{* Max Shifts: 12.00} \\
\texttt{* Min Shifts: 6.00} \\
\texttt{* End Distance Sum: 30.04} \\
\texttt{* Init Distance Sum: 30.04} \\
\texttt{* Expected Distance Sum: 29.33} \\
\texttt{* Penalties: 17.00} \\
\texttt{* Max Penalties: 16.00} \\
\texttt{* Rounds: 3.00} \\
\texttt{* Max Rounds: 28.00} \\
\texttt{* Object Count: 7.00} \\
            }
        }
    }
    & & \\ \\

\end{supertabular}
}

\end{document}
